\documentclass[a4paper,12pt]{ctexart}
\usepackage{../../teaching-style}

\begin{document}
\pagestyle{empty}

\maintitle{代数 上册}
\begin{center}
  \large 低学段 · 方程与不等式
\end{center}
\vspace{0.3cm}

\chaptertitle{第一章\quad 用字母表示数与简易方程}

\sectiontitle{1.1\quad 用字母表示数}

\knowbox{
用字母表示数可以简明地表达数量关系、运算律和公式。\\
常见用法：\\
\quad $\bullet$ 运算律：$a+b=b+a$，$a\times b=b\times a$\\
\quad $\bullet$ 面积公式：$S=ab$（长方形），$S=a^2$（正方形）\\
\quad $\bullet$ 数量关系：路程$=$速度$\times$时间，$s=vt$}

\begin{example}
用字母表示：
(1) $a$与$5$的和 \qquad (2) $x$的$3$倍减$y$的$\dfrac{1}{2}$\\
(3) $m$与$n$的积减去$m$与$n$的差
\end{example}
\begin{solution}
(1) $a+5$ \qquad (2) $3x-\dfrac{1}{2}y$ \qquad (3) $mn-(m-n)$
\end{solution}

\sectiontitle{1.2\quad 方程的概念与等式的性质}

\knowbox{
含有未知数的等式叫做方程。\\
\textbf{等式的性质}\\
性质1：等式两边同时加（或减）同一个数，等式仍然成立。\\
性质2：等式两边同时乘（或除以）同一个非零数，等式仍然成立。}

\begin{example}
利用等式的性质解方程：$3x+5=20$
\end{example}
\begin{solution}
两边减$5$：$3x=15$\\
两边除以$3$：$x=5$
\end{solution}

\sectiontitle{1.3\quad 解方程与验算}

\knowbox{
\textbf{移项}：把方程中的某项从一边移到另一边，要改变符号。\\
移项法则：$ax+b=c \rightarrow ax=c-b$\\
验算：把解代入原方程，检查左右两边是否相等。}

\begin{example}
解方程：$5x-7=3x+9$
\end{example}
\begin{solution}
移项：$5x-3x=9+7$\\
合并：$2x=16$\\
系数化$1$：$x=8$\\
验算：左边$=5\times8-7=33$，右边$=3\times8+9=33$，$x=8$是原方程的解。
\end{solution}

\sectiontitle{习题1}

\begin{exercise}
\item 用字母表示：
  \begin{subexercise}
    \item $x$的$2$倍与$y$的$3$倍的和
    \item $a$与$b$的差的平方
    \item $m$的$\dfrac{3}{4}$比$n$多$5$
    \item 三个连续整数，中间一个是$n$，另外两个是什么
  \end{subexercise}

\item 解方程：
  \begin{subexercise}
    \item $x+17=35$
    \item $x-12=28$
    \item $5x=75$
    \item $x\div6=12$
    \item $3x+7=28$
    \item $4x-9=23$
  \end{subexercise}

\item 解方程：
  \begin{subexercise}
    \item $2x+5=3x-1$
    \item $4x-3=2x+7$
    \item $5x+12=2x+30$
    \item $7x-8=3x+16$
    \item $6x-5=4x+3$
    \item $9x+4=5x+28$
  \end{subexercise}

\item 列方程解应用题：
  \begin{subexercise}
    \item 一个数的$3$倍加上$5$等于$26$，求这个数。
    \item 小明有$x$本书，借给同学$15$本后还剩$23$本，小明原有多少本？
    \item 长方形的周长为$36$厘米，宽为$6$厘米，长是多少厘米？
  \end{subexercise}
\end{exercise}

\newpage

\chaptertitle{第二章\quad 一元一次方程}

\sectiontitle{2.1\quad 去括号与去分母}

\knowbox{
\textbf{去括号法则}\\
$+(a-b)=a-b$，$-(a-b)=-a+b$\\
$a(b+c)=ab+ac$\\
\textbf{去分母}\\
方程两边同时乘以各分母的最小公倍数。}

\begin{example}
解方程：$2(3x-1)-3(2x+1)=4x-5$
\end{example}
\begin{solution}
去括号：$6x-2-6x-3=4x-5$\\
整理：$-5=4x-5$\\
移项：$-4x=0$\\
$x=0$
\end{solution}

\begin{example}
解方程：$\dfrac{2x-1}{3}-\dfrac{3x-4}{4}=1$
\end{example}
\begin{solution}
去分母（两边同乘$12$）：$4(2x-1)-3(3x-4)=12$\\
去括号：$8x-4-9x+12=12$\\
整理：$-x+8=12$\\
移项：$-x=4$\\
$x=-4$
\end{solution}

\sectiontitle{2.2\quad 一元一次方程的应用}

\knowbox{
列方程解应用题的步骤：\\
1. 审题，设未知数 \quad 2. 找等量关系\\
3. 列方程 \quad 4. 解方程 \quad 5. 检验并作答}

\begin{example}
甲、乙两地相距$300$千米。一辆汽车从甲地开往乙地，每小时行$60$千米；另一辆汽车从乙地开往甲地，每小时行$40$千米。两车同时出发，几小时后相遇？
\end{example}
\begin{solution}
设$x$小时后相遇。\\
等量关系：甲车路程$+$乙车路程$=300$\\
$60x+40x=300$\\
$100x=300$\\
$x=3$\\
答：$3$小时后相遇。
\end{solution}

\sectiontitle{习题2}

\begin{exercise}
\item 解方程：
  \begin{subexercise}
    \item $3(x+2)-2(x-1)=10$
    \item $4(2x-1)-3(x+2)=7$
    \item $5(2x-3)-2(3x-1)=5$
    \item $2(x-3)-3(x+1)=4(x-2)$
  \end{subexercise}

\item 解方程：
  \begin{subexercise}
    \item $\dfrac{x}{2}-\dfrac{x}{3}=5$
    \item $\dfrac{2x+1}{3}-\dfrac{x-1}{6}=2$
    \item $\dfrac{3x-2}{4}-\dfrac{2x+1}{3}=1$
    \item $\dfrac{x+4}{5}-\dfrac{x-1}{2}=1$
  \end{subexercise}

\item 列方程解应用题：
  \begin{subexercise}
    \item 一个数的$\dfrac{2}{3}$比它的$\dfrac{1}{2}$大$5$，求这个数。
    \item 把$140$本书分给两个班，甲班比乙班的$2$倍还多$20$本，两班各得多少本？
    \item 甲乙两人共存款$3600$元，甲存款是乙的$1.4$倍，甲乙各存款多少元？
    \item 甲比乙大$6$岁，$5$年前甲的年龄是乙的$2$倍，甲乙现在各几岁？
  \end{subexercise}
\end{exercise}

\newpage

\chaptertitle{第三章\quad 二元一次方程组}

\sectiontitle{3.1\quad 代入消元法}

\knowbox{
代入消元法步骤：\\
1. 把一个方程变形为用含一个未知数的式子表示另一个未知数。\\
2. 代入另一个方程，消去一个未知数。\\
3. 解一元一次方程。\\
4. 回代求另一个未知数。}

\begin{example}
解方程组：$\begin{cases}x+y=10\\2x-y=5\end{cases}$
\end{example}
\begin{solution}
由$\;$\;$(1)$得$y=10-x$，代入$(2)$：\\
$2x-(10-x)=5$\\
$2x-10+x=5$\\
$3x=15$，$x=5$\\
$y=10-5=5$\\
$\therefore\begin{cases}x=5\\y=5\end{cases}$
\end{solution}

\sectiontitle{3.2\quad 加减消元法}

\knowbox{
加减消元法步骤：\\
1. 将两个方程中同一个未知数的系数变成相等或互为相反数。\\
2. 两式相加或相减，消去一个未知数。\\
3. 解一元一次方程。\\
4. 回代求另一个未知数。}

\begin{example}
解方程组：$\begin{cases}3x+2y=13\\2x+3y=12\end{cases}$
\end{example}
\begin{solution}
$(1)\times2$：$6x+4y=26$\\
$(2)\times3$：$6x+9y=36$\\
相减：$-5y=-10$，$y=2$\\
代入$(1)$：$3x+4=13$，$x=3$\\
$\therefore\begin{cases}x=3\\y=2\end{cases}$
\end{solution}

\sectiontitle{习题3}

\begin{exercise}
\item 用代入消元法解：
  \begin{subexercise}
    \item $\begin{cases}x-y=3\\2x+y=12\end{cases}$
    \item $\begin{cases}x+2y=8\\3x-y=7\end{cases}$
    \item $\begin{cases}3x-2y=7\\x+y=4\end{cases}$
  \end{subexercise}

\item 用加减消元法解：
  \begin{subexercise}
    \item $\begin{cases}2x+3y=12\\3x+2y=13\end{cases}$
    \item $\begin{cases}5x+2y=25\\3x+4y=15\end{cases}$
    \item $\begin{cases}4x-3y=11\\3x+2y=4\end{cases}$
  \end{subexercise}

\item 列方程组解应用题：
  \begin{subexercise}
    \item 甲、乙两数的和为$32$，甲数的$3$倍等于乙数的$5$倍，求甲、乙两数。
    \item 篮球和排球共$12$个，总价$560$元。篮球每个$50$元，排球每个$40$元，篮球和排球各几个？
    \item 客车和货车共$12$辆，共载$420$人。客车每辆载$40$人，货车每辆载$30$人，客车和货车各几辆？
  \end{subexercise}
\end{exercise}

\newpage

\chaptertitle{第四章\quad 一元二次方程}

\sectiontitle{4.1\quad 一元二次方程的概念}

\knowbox{
形如$ax^2+bx+c=0(a\neq0)$的方程叫做一元二次方程。\\
$a$为二次项系数，$b$为一次项系数，$c$为常数项。}

\sectiontitle{4.2\quad 解法一：因式分解法}

\knowbox{
将方程化为$(x-m)(x-n)=0$的形式，则$x=m$或$x=n$。\\
常用方法：提公因式、十字相乘法。}

\begin{example}
解方程：$x^2-5x+6=0$
\end{example}
\begin{solution}
因式分解：$(x-2)(x-3)=0$\\
$\therefore x_1=2$，$x_2=3$
\end{solution}

\sectiontitle{4.3\quad 解法二：配方法}

\begin{example}
解方程：$x^2+6x-7=0$
\end{example}
\begin{solution}
移项：$x^2+6x=7$\\
配方：$x^2+6x+9=16$，即$(x+3)^2=16$\\
$x+3=\pm4$\\
$\therefore x_1=1$，$x_2=-7$
\end{solution}

\sectiontitle{4.4\quad 解法三：公式法}

\knowbox{
求根公式：$x=\dfrac{-b\pm\sqrt{b^2-4ac}}{2a}$\\
判别式$\Delta=b^2-4ac$：\\
$\Delta>0$，两个不等实根；\\
$\Delta=0$，两个相等实根；\\
$\Delta<0$，无实根。}

\begin{example}
解方程：$2x^2-5x+2=0$
\end{example}
\begin{solution}
$a=2$，$b=-5$，$c=2$\\
$\Delta=b^2-4ac=25-16=9>0$\\
$x=\dfrac{5\pm\sqrt9}{4}=\dfrac{5\pm3}{4}$\\
$\therefore x_1=2$，$x_2=\dfrac12$
\end{solution}

\sectiontitle{4.5\quad 根与系数的关系（韦达定理）}

\knowbox{
若$x_1$、$x_2$是$ax^2+bx+c=0(a\neq0)$的两个根，则\\
$x_1+x_2=-\dfrac{b}{a}$，\quad $x_1x_2=\dfrac{c}{a}$}

\sectiontitle{习题4}

\begin{exercise}
\item 用因式分解法解：
  \begin{subexercise}
    \item $x^2-9=0$
    \item $x^2-3x=0$
    \item $x^2+7x+12=0$
    \item $x^2-4x-5=0$
    \item $3x^2-5x+2=0$
    \item $2x^2+x-3=0$
  \end{subexercise}

\item 用配方法解：
  \begin{subexercise}
    \item $x^2+4x-5=0$
    \item $x^2-6x+8=0$
    \item $x^2-8x-9=0$
  \end{subexercise}

\item 用公式法解：
  \begin{subexercise}
    \item $x^2-5x+3=0$
    \item $2x^2-7x+3=0$
    \item $3x^2+5x-2=0$
  \end{subexercise}

\item 已知方程$x^2-5x+6=0$的两根为$x_1$、$x_2$，求：
  \begin{subexercise}
    \item $x_1+x_2$
    \item $x_1x_2$
    \item $\dfrac{1}{x_1}+\dfrac{1}{x_2}$
    \item $x_1^2+x_2^2$
  \end{subexercise}
\end{exercise}

\newpage

\chaptertitle{第五章\quad 不等式与不等式组}

\sectiontitle{5.1\quad 不等式的性质}

\knowbox{
性质1：若$a>b$，则$a\pm c>b\pm c$\\
性质2：若$a>b$，$c>0$，则$ac>bc$，$\dfrac{a}{c}>\dfrac{b}{c}$\\
性质3：若$a>b$，$c<0$，则$ac<bc$，$\dfrac{a}{c}<\dfrac{b}{c}$\\
注意：乘除负数时，不等号方向改变！}

\sectiontitle{5.2\quad 一元一次不等式的解法}

\begin{example}
解不等式：$2(x-3)\ge5x-6$
\end{example}
\begin{solution}
去括号：$2x-6\ge5x-6$\\
移项：$2x-5x\ge-6+6$\\
合并：$-3x\ge0$\\
系数化$1$：$x\le0$\\
解集在数轴上表示为：$\bullet\leftarrow$（原点处实心点）。
\end{solution}

\sectiontitle{5.3\quad 一元一次不等式组}

\knowbox{
大大取大，小小取小，大小小大中间找，大大小小找不到。\\
即：$x>a$且$x>b(a>b)\Rightarrow x>a$\\
$x<a$且$x<b(a<b)\Rightarrow x<a$\\
$x>a$且$x<b(a<b)\Rightarrow a<x<b$\\
$x<a$且$x>b(a<b)\Rightarrow$无解}

\sectiontitle{习题5}

\begin{exercise}
\item 解不等式：
  \begin{subexercise}
    \item $3x+5<2x+8$
    \item $4(x-2)>3(x-1)$
    \item $2(x+1)-3(x-2)\ge6$
    \item $\dfrac{x}{2}-\dfrac{x-1}{3}\ge1$
  \end{subexercise}

\item 解不等式组：
  \begin{subexercise}
    \item $\begin{cases}x+3>0\\2x-1<5\end{cases}$
    \item $\begin{cases}3x-2>4\\x+1<5\end{cases}$
    \item $\begin{cases}2x-1\ge x+1\\x-2\le 2x\end{cases}$
    \item $\begin{cases}x-2<0\\x+1>0\\2x-5\le1\end{cases}$
  \end{subexercise}

\item 含参讨论：$ax<b$的解集是什么？（提示：分$a>0$、$a<0$、$a=0$三种情况讨论）

\item 应用：某商店销售一种商品，进价为每件$40$元，售价不低于$50$元，且利润不超过进价的$30\%$，求售价的取值范围。
\end{exercise}

\newpage

\chaptertitle{第六章\quad 分式方程与根式方程}

\sectiontitle{6.1\quad 分式方程}

\knowbox{
分母中含未知数的方程叫分式方程。\\
解法：去分母（两边同乘最简公分母）$\rightarrow$解整式方程$\rightarrow$验根。\\
验根：使分母为零的根是增根，要舍去。}

\begin{example}
解方程：$\dfrac{2}{x-1}-\dfrac{3}{x}=1$
\end{example}
\begin{solution}
两边同乘$x(x-1)$：$2x-3(x-1)=x(x-1)$\\
整理：$2x-3x+3=x^2-x$，即$-x+3=x^2-x$\\
$x^2=3$，$x=\pm\sqrt3$\\
经检验，$x=\pm\sqrt3$都是原方程的解。
\end{solution}

\sectiontitle{6.2\quad 根式方程}

\knowbox{
根号内含有未知数的方程叫根式方程。\\
解法：移项使根号单独在一边$\rightarrow$两边平方$\rightarrow$解方程$\rightarrow$验根。\\
注意：平方后可能产生增根，必须验根。}

\begin{example}
解方程：$\sqrt{x+2}=x$
\end{example}
\begin{solution}
两边平方：$x+2=x^2$\\
整理：$x^2-x-2=0$，$(x-2)(x+1)=0$\\
$x=2$或$x=-1$\\
经检验，$x=2$是原方程的解，$x=-1$是增根（因为$\sqrt{x+2}=1\neq-1$）。
\end{solution}

\sectiontitle{习题6}

\begin{exercise}
\item 解分式方程：
  \begin{subexercise}
    \item $\dfrac{3}{x}=\dfrac{2}{x-1}$
    \item $\dfrac{2}{x+1}-\dfrac{1}{x}=1$
    \item $\dfrac{x}{x-1}-\dfrac{2}{x+1}=1$
    \item $\dfrac{1}{x-2}+3=\dfrac{2-x}{x-2}$
  \end{subexercise}

\item 解根式方程：
  \begin{subexercise}
    \item $\sqrt{x-1}=2$
    \item $\sqrt{2x+3}=x$
    \item $\sqrt{x+6}=x$
    \item $\sqrt{x+5}-\sqrt{x}=1$
  \end{subexercise}
\end{exercise}

\newpage

\chaptertitle{综合复习}

\begin{exercise}
\item 解方程：$3(2x-1)-2(3x-5)=7$
\item 解方程：$\dfrac{3x+1}{2}-\dfrac{2x-3}{3}=5$
\item 解方程组：$\begin{cases}2x+3y=8\\3x-2y=-1\end{cases}$
\item 解方程：$x^2-6x-7=0$
\item 解不等式组：$\begin{cases}2x-1>5\\3x-7\le8\end{cases}$
\item 解分式方程：$\dfrac{2}{x+1}+\dfrac{3}{x-1}=2$
\item 解根式方程：$\sqrt{2x+1}-\sqrt{x}=1$
\item 应用：一个两位数，十位数字与个位数字之和为$9$。将这个两位数对调后，新数比原数小$9$。求原数。
\item 应用：某种商品进价$80$元，标价$120$元。为了促销，商店决定打折销售但要保证利润率不低于$20\%$，则最多打几折？
\item 应用：甲乙两车分别从相距$360$千米的A、B两城同时出发相向而行，$2.4$小时后相遇。已知甲车速度是乙车的$1.2$倍，求两车速度。
\end{exercise}

\end{document}
